\clearpage
\begin{trapbox}
\begin{center}
\begin{tikzpicture}[font=\small,>=stealth]
\draw[gray!40,rounded corners] (0,0) rectangle (14.2,9);
\node[font=\small\bfseries] at (7.1,8.35) {防混淆：内分线沿切线；顶点处三角形退化};

\begin{scope}[shift={(3.6,4.7)},scale=.4]

\coordinate (L) at (-5,0); \coordinate (R) at (5,0);
\coordinate (P) at (3.75,3); \coordinate (M) at (2.4,0);
\draw[blue!75!black,thick] (L)--(P);
\draw[orange!85!black,thick] (P)--(R);
\draw[gray] (L)--(R);
\draw[magenta!80!black,very thick] (M)--(P);
\foreach \q in {L,R,P,M} \fill (\q) circle (3pt);
\node[below left] at (L) {$F_1$}; \node[below right] at (R) {$F_2$};
\node[above left] at (P) {$P$}; \node[below left] at (M) {$M$};

\draw[magenta!80!black] (1.7,-1.5556)--(4.9,5.5556);
\draw[gray,dashed] (0,4.6875)--(7,1.5375);
\node[above] at (6,2.1) {法线};
\end{scope}
\node at (3.7,2.5) {紫实线：内角平分线／切线};
\node at (3.7,1.6) {灰虚线：外角平分线／法线};
\begin{scope}[shift={(10.8,4.7)},scale=.42]
\draw[gray] (-5.8,0)--(6,0);
\draw[teal!75!black,domain=-.55:.55,samples=60,variable=\t] plot ({3*cosh(\t)},{4*sinh(\t)});
\draw[magenta!80!black] (3,-3)--(3,3);
\fill (-5,0) circle(4pt) node[below] {$F_1$};
\fill (5,0) circle(4pt) node[below] {$F_2$};
\fill (3,0) circle(4pt);\node[above left] at (3,0) {$P=A$};
\end{scope}
\node at (10.7,2.5) {$P,F_1,F_2$ 共线：没有三角形};
\node at (10.7,1.6) {顶点切线 $x=a$ 仍存在};
\node at (7.1,.65) {必须保留条件 $y_0\ne0$，才能谈焦点三角形的顶角内分线。};

\end{tikzpicture}
\end{center}

\begin{itemize}
\item \textbf{内角与外角混淆：}内角平分线沿切线，外角平分线沿法线；两者垂直，见上图。
\item \textbf{倒角公式漏写分母：}正确公式是 $\tan(\alpha-\beta)=(\tan\alpha-\tan\beta)/(1+\tan\alpha\tan\beta)$，不是两个斜率直接相减。
\item \textbf{正切相等不等于已经选对“内角”：}方法六会产生 $m=a^2/x_0$ 与 $m=c^2x_0/a^2$ 两根；后者满足 $|m|>c$，是外角平分线，不能误选。判据是内角交点 $-c<m<c$。
\item \textbf{斜率法的竖直线陷阱：}当 $x_0=\pm c$ 时，某条焦点连线的斜率不存在；若 $m=x_0$，$PM$ 的斜率也不存在。不能除以零；改用向量 $\det/\mathrm{dot}$ 写夹角正切（见方法六）。
\item \textbf{分母与约分陷阱：}推导前检查 $y_0\ne0$，正切差角的分母非零；先保留两个候选根，再用几何位置筛选，而不是提前丢掉一个根。
\item \textbf{比值方向颠倒：}必须同侧对应：$F_1M/MF_2=PF_1/PF_2$，见普通三角形配图。
\item \textbf{把左支的截距写正：}$x_0$ 与 $m$ 同号，统一用 $x_0m=a^2$；且 $0<|m|<a<c$，见例1、例2对称图。
\item \textbf{混淆切点与交点：}$T=(\operatorname{sgn}(x_0)a,0)$ 是圆的轴上切点，$M=(a^2/x_0,0)$ 是角平分线的轴上交点。非顶点时 $T\ne M$，见方法四图。
\item \textbf{漏掉非退化条件：}必须 $y_0\ne0$。顶点处不能套三角形角平分线定理；顶点切线仍存在，见上图。
\item \textbf{循环论证：}方法三引用已知切线性质，只适合快速求截距；独立证明应走方法一，再由 $P,M$ 写出切线方程。
\item \textbf{坐标或焦距关系写错：}$P=(x_0,y_0)$、$M=(m,0)$ 满足乘积关系 $x_0m=a^2$；双曲线用 $c^2=a^2+b^2$。
\end{itemize}
\end{trapbox}
